\section{Experimental Details}
\label{app:details}

\subsection{Hardware and Software}

All perplexity measurements were conducted on NVIDIA RTX A6000 GPUs (49GB VRAM) with CUDA~12.5, Python~3.12.8, PyTorch~2.5.1+cu124, and Hugging Face Transformers~5.5.0. Scenario and variant generation used the OpenAI API (\gptfour) with estimated cost of \$3--5.

\subsection{Scenario Domains}

\Tabref{tab:domains} lists example scenarios from each domain.

\begin{table}[h]
    \centering
    \caption{Example scenarios from each communicative domain.}
    \begin{tabular}{@{}lp{10cm}@{}}
        \toprule
        \textbf{Domain} & \textbf{Example Scenario} \\
        \midrule
        Workplace & ``Write an email to your manager explaining that you will miss Friday's deadline due to a family emergency.'' \\
        Personal & ``Write a message to your neighbor asking them to keep the noise down during weeknight evenings.'' \\
        Service & ``Write a complaint to customer support about a defective product you received last week.'' \\
        \bottomrule
    \end{tabular}
    \label{tab:domains}
\end{table}

\subsection{Variant Generation Prompts}

For each axis, we instructed \gptfour to generate two response variants with the following template (abridged):

\begin{quote}
\small
\texttt{Given the scenario below, write two responses to the user's request. Response~A should be [Pole~A description]. Response~B should be [Pole~B description]. Both responses should: (1) address the same scenario, (2) convey the same core message, (3) be similar in length ($\pm$20\%), and (4) differ primarily on the [axis name] dimension.}
\end{quote}

\subsection{Reproducibility}

All experiments use NumPy and PyTorch random seed 42. \gptfour was called with temperature 0.7--0.8 for generation and 0.3 for introspection. Complete raw outputs, including all scenario texts, variant pairs, log-likelihood scores, and introspection responses, are available in the \texttt{results/} directory of our code repository.
